\documentclass[12pt]{article}
\usepackage[T1]{fontenc}
\usepackage{amsmath,amssymb,amsthm}
\usepackage{newtxtext,newtxmath}
\usepackage{array}
\usepackage{url}
\newtheorem{theorem}{Theorem}[section]
\newtheorem{corollary}[theorem]{Corollary}
\newtheorem{conjecture}[theorem]{Conjecture}
\newtheorem{lemma}[theorem]{Lemma}
\newtheorem{proposition}[theorem]{Proposition}
\newtheorem{definition}[theorem]{Definition}
\newtheorem{example}[theorem]{Example}
\newtheorem{remark}[theorem]{Remark}
\newtheorem{problem}[theorem]{Problem}
\numberwithin{equation}{section}
\newcommand{\F}{\mathbb{F}}
\newcommand{\PP}{\mathbb{P}}
\newcommand{\A}{\mathbb{A}}
\newcommand{\supp}{\operatorname{supp}}
\newcommand{\Aut}{\operatorname{Aut}}
\newcommand{\Stab}{\operatorname{Stab}}
\newcommand{\ev}{\operatorname{ev}}
\newcommand{\wt}{\operatorname{wt}}
\newcommand{\RS}{\operatorname{RS}}
\newcommand{\PAut}{\operatorname{PAut}}
\newcommand{\AGL}{\operatorname{AGL}}
\newcommand{\PSL}{\operatorname{PSL}}
\newcommand{\PGL}{\operatorname{PGL}}
\newcommand{\Span}{\operatorname{span}}
\begin{document}
\title{Permutation decoding of codes from $y^2=x^p-x$:\\
corrected results, counterexamples, and partial decoding}
\author{David Joyner\thanks{Original author; Mathematics Department,
U.S. Naval Academy, Annapolis, MD 21402; wdj@usna.edu.}
\and ChatGPT (OpenAI)\thanks{AI-assisted mathematical revision,
proof checking, and exposition; credited as a co-author at the
request of the original author.}}
\date{November 2026 --- mathematical revision, Step 1}
\maketitle
\begin{abstract}
This is an AI-aided paper, revising the original note on permutation decoding
of algebraic-geometric codes associated with the hyperelliptic curves
$y^2=x^p-x$.  The length-$p$ codes over $\F_p$ are identified as ordinary
Reed--Solomon codes.  A general counting obstruction is used to disprove
two proposed full permutation-decoding assertions, with explicit examples
at $p=37$ and $p=7$.  On the positive side, the full affine permutation group
gives a two-error permutation-decoding set whenever the corresponding
Reed--Solomon code corrects at least two errors.  We also correct the
stabilizer and representation-theoretic assertions and exhibit a
six-dimensional Riemann--Roch module which is an invariant submodule of a
permutation module but is not itself a permutation module.  We formulate
a more precise partial-decoding problem for the codes over $\F_{p^2}$.
\end{abstract}

\section{Introduction}
Let $X$ be a smooth, projective, geometrically integral curve defined over a
finite field $F=\F_q$.  Let $D$ be a divisor on $X$ defined over $F$ and let
\[
E=P_1+\cdots+P_n
\]
be a sum of distinct $F$-rational points with
$\supp(D)\cap\supp(E)=\varnothing$.  The Riemann--Roch space is
\[
L(D)=\{f\in F(X)^\times:(f)+D\geq0\}\cup\{0\}.
\]
The associated algebraic-geometric (AG) evaluation code is
\begin{equation}\label{eq:AGcode}
C_X(D,E)=\{(f(P_1),\ldots,f(P_n)):f\in L(D)\}\subseteq F^n.
\end{equation}
Write $\ev_E:L(D)\to F^n$ for the evaluation map.  A standard
sufficient condition for its injectivity is $\deg D<n$: a nonzero
function in its kernel would have at least $n$ zeros outside
$\supp(D)$, contrary to $(f)+D\geq0$.

If a finite group $G\leq\Aut_F(X)$ preserves the divisors $D$ and $E$,
then $G$ permutes the evaluation coordinates and preserves the code.
This provides an opportunity to exploit geometric symmetry in decoding.
The distinction between the action on coordinates, the induced action on
codewords, and the representation on $L(D)$ is important and will be made
precise in Section~\ref{sec:modules}.

The original note proposed efficient permutation decoding for two families
of codes arising from
\begin{equation}\label{eq:curve}
X:\qquad y^2=x^p-x.
\end{equation}
Both proposed full-decoding statements require qualification.  We retain
the curve-theoretic setup and the original computational illustration but
replace the conjectured universal assertions by provable statements,
explicit obstructions, and a narrower open problem.

Some of the original ideas arose in joint research with Amy Ksir and
Will Traves; see \cite{JK1,JK2,JT}.  The present revision keeps that
historical context while distinguishing the mathematical claims
proved here from the original conjectures.

\section{The curve and its finite group actions}\label{sec:curve}
Throughout this section $p\geq5$ is an odd prime, $g=(p-1)/2$,
and $X$ denotes the smooth projective model of \eqref{eq:curve}.
The condition $p\geq5$ avoids the exceptional genus-one case $p=3$.
The polynomial $x^p-x$ is squarefree, so $X$ is hyperelliptic of genus $g$.
There is a unique point $P_\infty$ above infinity under the
hyperelliptic projection $x:X\to\PP^1$.
One should not confuse $P_\infty$ with a finite point $(a,0)$;
in the ordinary plane projective closure the point at infinity is
$[0:1:0]$, and the smooth projective model is used for divisor
computations.

\begin{lemma}\label{lem:poles}
The valuations at the point $P_\infty$ satisfy
\[
\operatorname{ord}_{P_\infty}(x)=-2,\qquad
\operatorname{ord}_{P_\infty}(y)=-p.
\]
For every integer $m\geq0$ the following is an $F$-basis of
$L(mP_\infty)$:
\begin{equation}\label{eq:monomial-basis}
\{x^i:i\geq0,\ 2i\leq m\}\ \cup\
\{x^iy:i\geq0,\ 2i+p\leq m\}.
\end{equation}
\end{lemma}
\begin{proof}
The valuations follow from the degree-two map $x$ and the equation
$y^2=x^p-x$.  Every function regular away from $P_\infty$ is
represented by an element of $F[x,y]/(y^2-x^p+x)$, which has a unique
expression $A(x)+yB(x)$.  The pole orders of the monomials in
\eqref{eq:monomial-basis} are distinct: the first are even and the
second are odd.  Therefore leading pole terms cannot cancel, and
exactly the indicated monomials belong to $L(mP_\infty)$.
\end{proof}

\begin{definition}
A finite subgroup $G\leq\Aut_F(X)$ will be called \emph{large}
(over $F$) if $|G|>|X(F)|$.
\end{definition}
\begin{lemma}\label{lem:large}
If $G$ is large, every $F$-rational point is ramified in
the quotient map $X\to X/G$.
\end{lemma}
\begin{proof}
If $P\in X(F)$ had trivial stabilizer, the orbit--stabilizer theorem
would give $|G\cdot P|=|G|>|X(F)|$, impossible.
For the quotient by a finite automorphism group, a nontrivial
geometric stabilizer is the ramification group at that point.
\end{proof}

\subsection{Automorphisms over $\F_p$}\label{subsec:fp}
The $\F_p$-rational points are
\begin{equation}\label{eq:fp-points}
X(\F_p)=\{P_\infty\}\cup\{P_a=(a,0):a\in\F_p\},
\qquad |X(\F_p)|=p+1.
\end{equation}
The hyperelliptic involution is $\iota(x,y)=(x,-y)$.
The branch locus of $x:X\to\PP^1$ is $\PP^1(\F_p)$.
A M\"obius transformation
\[
x\longmapsto\frac{ax+b}{cx+d},\qquad
M=\begin{pmatrix}a&b\\c&d\end{pmatrix}\in\operatorname{GL}_2(\F_p),
\]
admits a lift to the curve of the form
\begin{equation}\label{eq:lift}
(x,y)\longmapsto
\left(\frac{ax+b}{cx+d},\frac{\eta y}{(cx+d)^{(p+1)/2}}\right),
\qquad \eta^2=\det M.
\end{equation}
Indeed $x'^p-x'=(\det M)(x^p-x)/(cx+d)^{p+1}$.
Over $\F_p$ such a lift exists when $\det M$ is a square,
and each allowable projectivity has the two lifts differing by $\iota$.
Thus the resulting geometric group $G_p$ maps onto
$\PSL_2(p)$ with kernel $\langle\iota\rangle$, and
\begin{equation}\label{eq:gp-order}
|G_p|=p(p^2-1).
\end{equation}
For the present genus-$g\geq2$ curve these lifts exhaust the
$\F_p$-automorphisms, since every automorphism descends to a
projectivity preserving the hyperelliptic branch locus
(compare \cite{G}).

Explicit lifts include
\begin{equation}\label{eq:generators}
\begin{aligned}
\iota &: (x,y)\mapsto(x,-y),\\
\delta_a &: (x,y)\mapsto(a^2x,ay),\\
\tau_b &: (x,y)\mapsto(x+b,y),\\
\omega &: (x,y)\mapsto
\bigl(-1/x,\,y/x^{(p+1)/2}\bigr),
\end{aligned}
\end{equation}
where $a\in\F_p^\times$ and $b\in\F_p$.
The group $G_p$ is transitive on the $p+1$ points in
\eqref{eq:fp-points}.  The stabilizer
$H_p=\Stab_{G_p}(P_\infty)$ is generated by
$\iota$, the $\delta_a$, and the $\tau_b$, and has order
\begin{equation}\label{eq:Hp-order}
|H_p|=p(p-1).
\end{equation}
It is a solvable group.  On the $p$ finite branch points,
its induced coordinate-permutation group is
\begin{equation}\label{eq:square-affine}
\{x\mapsto ux+b:\ u\in(\F_p^\times)^2,\ b\in\F_p\},
\qquad\text{of order }p(p-1)/2,
\end{equation}
because the involution fixes every branch point.  This is generally
smaller than the \emph{full} permutation automorphism group of the
Reed--Solomon code encountered below.

\begin{proposition}\label{prop:H-comp}
For every $m\geq0$, the $H_p$-module $L(mP_\infty)$ has a
composition series with one-dimensional factors.  Equivalently,
its semisimplification is a direct sum of one-dimensional
$H_p$-modules (not, in general, $G_p$-modules).
\end{proposition}
\begin{proof}
Order the monomials in \eqref{eq:monomial-basis} by increasing
pole order.  The subspace of functions with pole order at most
$r$ is $H_p$-stable because $H_p$ fixes $P_\infty$.
Each successive nonzero quotient of this filtration is
one-dimensional: the pole orders are distinct.  Explicitly,
$\tau_b(x^iy^\epsilon)=(x+b)^iy^\epsilon$ is triangular,
$\delta_a(x^iy^\epsilon)=a^{2i+\epsilon}x^iy^\epsilon$,
and $\iota(x^iy^\epsilon)=(-1)^\epsilon x^iy^\epsilon$
for $\epsilon\in\{0,1\}$.  Thus the successive factors are
characters of $H_p$.  Notice that $\iota$ does \emph{not}
act trivially on terms involving $y$.
\end{proof}

\subsection{Automorphisms and rational orbits over $\F_{p^2}$}
\label{subsec:fp2}
Let $F=\F_{p^2}$.  Since every element of $\F_p^\times$
has a square root in $F$, the transformations
\eqref{eq:lift} lift all of $\PGL_2(p)$.  Consequently
\begin{equation}\label{eq:gq-order}
1\longrightarrow\langle\iota\rangle\longrightarrow
G_{p^2}\longrightarrow\PGL_2(p)\longrightarrow1,
\qquad |G_{p^2}|=2p(p^2-1).
\end{equation}
For $p\geq5$ this is the $F$-rational automorphism group
of $X$.  The stabilizer $H=\Stab_{G_{p^2}}(P_\infty)$
has order $2p(p-1)$; it acts by
\begin{equation}\label{eq:H-affine}
(x,y)\longmapsto (ux+b,\,\eta y),
\qquad u\in\F_p^\times,\quad b\in\F_p,\quad \eta^2=u.
\end{equation}

\begin{proposition}\label{prop:orbits}
If $p\equiv3\pmod4$, the $G_{p^2}$-orbits on $X(\F_{p^2})$ are
\[
O_1=X(\F_p),\quad |O_1|=p+1;
\qquad O_2=X(\F_{p^2})\setminus X(\F_p),
\quad |O_2|=2p(p-1).
\]
The subgroup $H$ in \eqref{eq:H-affine} acts regularly
(freely and transitively) on $O_2$.  Therefore for
$P_2\in O_2$ its stabilizer $H_2$ in $G_{p^2}$ has
order $p+1$, and $H\cap H_2=\{1\}$.
If $p\equiv1\pmod4$, then $X(\F_{p^2})=X(\F_p)$
has $p+1$ elements and is a single $G_{p^2}$-orbit.
\end{proposition}
\begin{proof}
For $x\in\F_p$ the equation $y^2=x^p-x$ gives $y=0$.
For $x\in\F_{p^2}\setminus\F_p$, put $z=x^p-x\ne0$.
Then $z^p=-z$ and
\[
z^{(p^2-1)/2}=(z^{p-1})^{(p+1)/2}
=(-1)^{(p+1)/2}.
\]
Thus $z$ is a square in $\F_{p^2}$ if and only if
$p\equiv3\pmod4$.  In that case each of the $p^2-p$
values $x\notin\F_p$ has two points above it, giving the
asserted orbit cardinality.

Fix a basis $1,\theta$ of $\F_{p^2}$ over $\F_p$.
Every $x\notin\F_p$ is uniquely $a+b\theta$ with
$a\in\F_p$, $b\in\F_p^\times$.  The affine group
$\AGL(1,p)$ acts freely and transitively on these $x$ values.
For each such affine map there are two lifts \eqref{eq:H-affine},
which interchange the two possible $y$ values.
Hence $H$ is transitive on $O_2$ and, as its order is
$|O_2|$, the action is regular.  The orbit $O_1$ is
transitive under the projective transformations already described.
The stabilizer assertions follow from orbit--stabilizer.
\end{proof}

For example, with $p\equiv3\pmod4$ the invariant divisor
$D_2=\sum_{P\in O_2}P$ satisfies
$\deg D_2=2p(p-1)>2g-2$, and Riemann--Roch gives
\[
\dim_F L(D_2)=2p(p-1)+1-g.
\]
This is part of the motivation for studying the larger
$\F_{p^2}$-orbit codes.

\section{AG codes and the relevant group representations}
\label{sec:modules}
Let $G\leq\Aut_F(X)$ be a finite group and suppose $D$ and $E$
are both $G$-invariant and have disjoint support.  Then $G$ acts
on $L(D)$ by pullback,
\[
(g\cdot f)(P)=f(g^{-1}P),
\]
and on the $n$ coordinate positions by permuting $\supp(E)$.
These actions are compatible with evaluation:
\begin{equation}\label{eq:equivariance}
\ev_E(g\cdot f)=\pi_E(g)\cdot\ev_E(f).
\end{equation}
The vector on the right is obtained by permuting the coordinates
of $\ev_E(f)$.  Thus the image of the positional action lies in
$\PAut(C_X(D,E))\leq S_n$.

\begin{proposition}\label{prop:actions}
Write $C=C_X(D,E)$.  There are distinct homomorphisms
\[
\pi_E:G\longrightarrow\PAut(C)\leq S_n,
\qquad \rho_C:G\longrightarrow\operatorname{GL}(C).
\]
The kernel of $\pi_E$ is exactly the pointwise stabilizer of
$\supp(E)$, without any hypothesis on $L(D)$.
If $L(D)$ separates the points in $\supp(E)$, then
$\ker\rho_C=\ker\pi_E$.
If $\ev_E$ is injective (for example, if $\deg D<n$),
then $\ev_E$ is an isomorphism of $F[G]$-modules from
$L(D)$ onto $C$; in particular $C$ is a submodule of the
\emph{ambient permutation module} $F^n$.
\end{proposition}
\begin{proof}
The first assertion about $\pi_E$ is the definition of a
faithful permutation action on the indicated coordinate set.
If $g\in\ker\rho_C$, then $f(P)=f(g^{-1}P)$ for
all $P\in\supp(E)$ and all $f\in L(D)$.  Separation of
these points forces $g^{-1}P=P$ for every $P$.
The reverse implication is immediate.  Finally, the evaluation
map is surjective onto $C$ by definition, and
\eqref{eq:equivariance} makes it equivariant; injectivity
therefore gives the required module isomorphism.
\end{proof}

\begin{remark}\label{rem:separation}
The assertion ``$D$ is ample, so $L(D)$ separates points''
requires correction.  A positive-degree divisor on a smooth
projective curve is ample, but its complete linear system need
not separate points.  A standard sufficient condition is
$\deg D\geq2g+1$: then $|D|$ is very ample and separates
points (where evaluation makes sense).  For example,
$D=P_\infty$ has positive degree, yet for $g\geq1$,
$L(P_\infty)=F$ and separates no pair of distinct points.
Also note that neither injectivity of evaluation nor point
separation is required for the coordinate action $\pi_E$.
\end{remark}

\subsection{Invariant submodules need not be permutation modules}
The original manuscript asserted that, under degree hypotheses
such as $n>\max(2\deg D,2g+2)$, the representation on $L(D)$
can be realized by permutation matrices.  This conclusion does
not follow from the equivariant embedding in
Proposition~\ref{prop:actions}.  The correct conclusion is
\emph{an invariant submodule of a permutation module}, not
necessarily a permutation module in its own right.

\begin{example}[A counterexample satisfying the original degree bounds]
\label{ex:not-permutation}
Let $p=7$, $F=\F_{49}$, $g=3$, and let $O_1,O_2$ be the
$G_{49}$-orbits from Proposition~\ref{prop:orbits}.  Put
\[
D=\sum_{P\in O_1}P,
\qquad E=\sum_{P\in O_2}P.
\]
These are $G_{49}$-invariant and disjoint, with
\[
\deg D=8,\quad n=\deg E=84,
\quad 84>\max\{2\cdot8,\,2\cdot3+2\}=16.
\]
Riemann--Roch gives
$\dim_F L(D)=8-3+1=6$, and $\ev_E$ is injective.
The functions $1$, $1/y$, and $x/y$ belong to $L(D)$:
$y$ has simple zeros at the seven finite branch points
and a pole of order seven at infinity, whereas $x$ has a
pole of order two at infinity.
The automorphism $\tau_1:(x,y)\mapsto(x+1,y)$ has order
seven and acts nontrivially, since
\[
\tau_1^*(x/y)=\frac{x+1}{y}=x/y+1/y.
\]
Any permutation matrix of size six having order dividing
seven is the identity, since a nontrivial cycle of length
seven cannot act on six letters.  Therefore the nontrivial
action of $\tau_1$ on $L(D)$ cannot arise from a permutation
representation of degree six.  This is a counterexample to
the original assertion, even though all its stated numerical
hypotheses hold.
\end{example}

\begin{remark}
Example~\ref{ex:not-permutation} shows why the word
\emph{submodule} in Proposition~\ref{prop:actions} cannot be
replaced by \emph{permutation module}.  The injectivity
hypothesis is needed for the asserted module isomorphism:
if it is dropped, $C$ is merely the quotient
$L(D)/\ker\ev_E$.  The indecomposable decomposition of the
resulting modular representations remains an interesting
question; see \cite{JK2,JT,S}.
\end{remark}

\section{Permutation decoding and a counting obstruction}
\label{sec:PD}
Let $C$ be an $[n,k,d]$ linear code over $F$ and put
\[
t=\left\lfloor\frac{d-1}{2}\right\rfloor,
\qquad r=n-k.
\]
Let $I$ be an information set (a set of $k$ coordinates on
which restriction from $C$ is an isomorphism), and set
$J=\{1,\dots,n\}\setminus I$ for the $r$ check positions.
After a suitable coordinate ordering and row operations the
generator and parity-check matrices have systematic forms
\[
G_0=(I_k\mid A),\qquad H_0=(-A^{\mathsf T}\mid I_r).
\]

\begin{lemma}[Syndrome criterion]\label{lem:syndrome}
Suppose $v=c+e$ with $c\in C$, $\wt(e)\leq t$.
The information coordinates of $v$ agree with those of $c$
if and only if $\wt(H_0v^{\mathsf T})\leq t$.
\end{lemma}
\begin{proof}
If the information coordinates are correct, $e$ is supported
in $J$, and $H_0e^{\mathsf T}$ equals its $r$ check entries;
its weight is at most $t$.  Conversely, write the information
part of $e$ as $u$ and put $s=H_0e^{\mathsf T}$.
The vector $(u,A^{\mathsf T}u)$ belongs to $C$.  If $u\ne0$, then
\[
(u,A^{\mathsf T}u)=e-(0,s),
\]
so it is a nonzero codeword of weight at most
$\wt(e)+\wt(s)\leq2t<d$, a contradiction.
Thus $u=0$.
\end{proof}

\begin{definition}\label{def:pd-set}
For $0\leq s\leq t$, a set
$S\subseteq\PAut(C)$ is an \emph{$s$-permutation-decoding set}
(or $s$-PD-set) relative to $I$ if, for every set $T$ of at
most $s$ coordinate positions, some $\sigma\in S$ sends
$T$ into $J$.  A $t$-PD-set is called a full PD-set.
\end{definition}

With an $s$-PD-set, permutation decoding tests $\sigma v$
for $\sigma\in S$ until the syndrome has weight at most $t$,
encodes the resulting information symbols, and applies
$\sigma^{-1}$ to the codeword.  Lemma~\ref{lem:syndrome}
justifies the algorithm for $\wt(e)\leq s$.
One may use the entire coordinate automorphism group as a
candidate set, but the existence of a full PD-set is not
a consequence of transitivity alone.

\begin{theorem}[Necessary counting bound]\label{thm:counting}
Let $C$ be an $[n,k,d]$ code and let $S$ be an $s$-PD-set
for an information set $I$, where $1\leq s\leq t$.
Then
\begin{equation}\label{eq:counting}
|S|\binom{n-k}{s}\ \geq\ \binom ns.
\end{equation}
Consequently, if
\begin{equation}\label{eq:counting-fails}
|\PAut(C)|\binom{n-k}{s}<\binom ns,
\end{equation}
no $s$-PD-set exists for any information set.
More generally, the same obstruction holds for any
specified subgroup $K\leq\PAut(C)$ upon replacing
$|\PAut(C)|$ by $|K|$.
\end{theorem}
\begin{proof}
For a fixed $\sigma$, there are exactly
$\binom{r}{s}$ subsets $T$ of size $s$ such that
$\sigma(T)\subseteq J$, namely the $s$-subsets of
$\sigma^{-1}(J)$.  The sets handled by the elements of $S$
must cover all $\binom ns$ error supports of size $s$.
Counting with possible repetitions proves
\eqref{eq:counting}; its contrapositive gives the rest.
\end{proof}

\begin{remark}
Inequality \eqref{eq:counting} is only a \emph{necessary}
condition: satisfying it does not guarantee the existence
of an $s$-PD-set.  Sharper lower bounds based on iterative
covering arguments are known in the permutation-decoding
literature \cite{HP}.
\end{remark}

\section{Codes over $\F_p$: the Reed--Solomon identification}
\label{sec:RS}
The original proposed conjecture for one-point codes of
length $p$ can be made more precise by identifying the codes
explicitly.  Here $E=\sum_{a\in\F_p}P_a$ and
$D=mP_\infty$ with $m\geq0$.

\begin{theorem}\label{thm:RS}
The AG evaluation code is exactly the full-length
Reed--Solomon code
\begin{equation}\label{eq:RS}
\begin{aligned}
C_X(mP_\infty,E)
 &=\{(f(a))_{a\in\F_p}: f\in\F_p[x],\ \deg f<k\},\\
k&=\min\left(p,\left\lfloor\frac m2\right\rfloor+1\right).
\end{aligned}
\end{equation}
For $1\leq k\leq p$, its parameters are
\[
[n,k,d]=[p,k,p-k+1].
\]
For $2\leq k\leq p-2$, its full permutation automorphism
group is $\AGL(1,p)$, of order $p(p-1)$.
\end{theorem}
\begin{proof}
At every evaluation point $P_a=(a,0)$, all monomials involving
$y$ vanish.  Lemma~\ref{lem:poles} therefore identifies the
image of evaluation with that of the polynomials
$1,x,\ldots,x^{\lfloor m/2\rfloor}$ at every element of
$\F_p$.  Reduction modulo $x^p-x$ shows that the image
has dimension $k$ as stated.  The Reed--Solomon distance
formula follows from the bound on the number of roots of a
nonzero polynomial of degree less than $k$, with equality
obtained by taking a product of $k-1$ distinct linear factors.
The determination of the full permutation group in the
nonexceptional dimension range is the classical full-length
Reed--Solomon result; see \cite{KW,CLLS}.  Each affine map
$x\mapsto ux+b$, $u\in\F_p^\times$, certainly preserves
the code by substitution, while these references show
that there are no other coordinate permutations in this range.
\end{proof}

\begin{remark}\label{rem:RSkernel}
When $m\geq p$, the function $y$ lies in $L(mP_\infty)$
and evaluates to zero at all the $P_a$.  Thus the code is
Reed--Solomon even though the evaluation map from the
larger Riemann--Roch space need not be injective.
The full affine group $\AGL(1,p)$ of the code is larger
than the geometric stabilizer image \eqref{eq:square-affine}:
a nonsquare multiplier generally cannot be lifted to an
$\F_p$-automorphism of the curve fixing $P_\infty$.
\end{remark}

\begin{theorem}[A positive partial-decoding result]
\label{thm:two-pd}
For $2\leq k\leq p-4$ the full group
$\AGL(1,p)$ is a $2$-PD-set for the code in
Theorem~\ref{thm:RS}, relative to \emph{every} information
set.  More generally, for any full-length Reed--Solomon
code with at least two check coordinates, every error
support of size at most two can be moved into the check
coordinates by an affine permutation.
\end{theorem}
\begin{proof}
The affine group $\AGL(1,p)$ is sharply two-transitive
on $\F_p$: for any distinct $a_1,a_2$ and distinct
$b_1,b_2$ there is a unique affine map taking
$a_i$ to $b_i$ for $i=1,2$.  Given an error support
of size two, choose any two distinct check positions and
apply the corresponding affine map.  The assertion for a
single error follows from transitivity, and the zero-error
case is immediate.  In the indicated range of $k$,
$d=p-k+1\geq5$, so $t\geq2$ and the construction is
indeed a $2$-PD-set in the standard decoding sense.
\end{proof}

\begin{example}[The original length-seven computation]
\label{ex:7-3-5}
For $p=7$ and $D=4P_\infty$ the basis of
$L(D)$ is $\{1,x,x^2\}$, so
$C_X(D,E)=\RS(\F_7,3)$ has parameters $[7,3,5]$.
The original GUAVA calculation gave the systematic matrix
\[
G_0=\begin{pmatrix}
1&0&0&6&4&6&4\\
0&1&0&6&3&3&1\\
0&0&1&3&1&6&3
\end{pmatrix},
\]
with permutation generators
$ (1,5,4)(2,3,7)$ and $(1,3,2,6,7,4)$.
These preserve the row space and generate a group of
order $42=7(7-1)$, namely $\AGL(1,7)$.
By Theorem~\ref{thm:two-pd} the group is a full
$2$-PD-set (because the decoding radius is $t=2$).
For each unordered pair of error positions, exactly
$(7-3)(7-3-1)=12$ affine permutations move both errors
into the four check positions; this count follows by
sharp two-transitivity on ordered pairs.
\end{example}


\subsection{Historical GUAVA transcript}
\label{subsec:guava}
For reproducibility we retain the short commands from the original note.
They serve as a record of how the $[7,3,5]$ example was investigated,
not as a proof of the universal conjecture.  The original typo
\texttt{Length(P)} has been corrected to \texttt{Length(Pts)}.
\begin{verbatim}
p:=7;
F:=GF(p);
R := PolynomialRing(F,2);
indets := IndeterminatesOfPolynomialRing(R);
x:=indets[1]; y:=indets[2];
Pts:=AffinePointsOnCurve(y^2-x^p+x,R,F);;
Length(Pts);
C:=OnePointAGCode(y^2-x^p+x,Pts,4,R);
MinimumDistance(C);
G:=GeneratorMat(C);
PutStandardForm(G);
Display(G);
P:=PermutationAutomorphismGroup(C);
IdGroup(P);
Wt2:=ConstantWeightSubcode(WholeSpaceCode(p,GF(p)),2);
for c in C do
 e:=Random(Wt2);; v:=PermutationDecode(C,c+e);
 if c<>v then
  Print("\\n Decoding fails. c,v = ",c," ",v,"\\n");
 fi;
od;
\end{verbatim}
These historical commands are version-dependent.  The mathematical
claims above do not rely on the random experiments.

\subsection{A counterexample to unrestricted permutation decoding}
\begin{theorem}\label{thm:counter-RS}
For $p=37$, $D=32P_\infty$ and
$E=\sum_{a\in\F_{37}}P_a$, the full permutation
automorphism group is \emph{not} a full PD-set for
$C_X(D,E)$, regardless of the information set.
Consequently the first universal full-decoding conjecture
of the original note is false.
\end{theorem}
\begin{proof}
Theorem~\ref{thm:RS} gives the Reed--Solomon parameters
\[
[n,k,d]=[37,17,21],\qquad r=20,\quad t=10,
\]
and the full group has order $37\cdot36=1332$.
But
\begin{align*}
1332\binom{20}{10}
 &=1332\cdot184756=246094992\\
 &<348330136=\binom{37}{10}.
\end{align*}
This contradicts the necessary bound
\eqref{eq:counting} for $s=t=10$.  Hence the
full automorphism group cannot be a $10$-PD-set.
\end{proof}

\section{Codes over $\F_{p^2}$: limits of the geometric stabilizer}
\label{sec:fp2-codes}
Assume $p\equiv3\pmod4$, $p\geq7$, and put
$E_2=\sum_{P\in O_2}P$, with $O_2$ from
Proposition~\ref{prop:orbits}.  The divisor
$D=mP_\infty$ is invariant under the point stabilizer
$H\leq G_{p^2}$, whose order is $2p(p-1)$.
The group $H$ acts regularly on the $n=2p(p-1)$
positions of the code $C_X(D,E_2)$.
This is a natural candidate for a geometric PD-set,
but it is not always large enough.

\begin{theorem}\label{thm:counter-fp2}
Let $p=7$, $F=\F_{49}$, $D=44P_\infty$ and
$E_2=\sum_{P\in O_2}P$.  The associated AG code has
length $84$, dimension $42$, and minimum distance at least
$40$.  The image of $H=\Stab_{G_{49}}(P_\infty)$,
which has order $84$, is \emph{not} a $7$-PD-set for
any information set.  Thus the second proposed
universal full-decoding conjecture is false for its
specified geometric stabilizer.
\end{theorem}
\begin{proof}
Here $g=3$, $\deg D=44<84=\deg E_2$, so evaluation
is injective.  Since $44>2g-2$, Riemann--Roch gives
\[
k=\ell(D)=\deg D+1-g=42.
\]
The usual AG bound gives $d\geq n-\deg D=40$,
so $t\geq\lfloor(40-1)/2\rfloor=19$.
In particular, seven-error PD is within the guaranteed
error-correction radius.  However
\begin{align*}
|H|\binom{n-k}{7}
&=84\binom{42}{7}=2266179552\\
&<4529365776=\binom{84}{7}.
\end{align*}
Proposition~\ref{prop:orbits} shows that $H$ acts
regularly on $O_2$, hence its positional image has
exactly $84$ elements.  Theorem~\ref{thm:counting}
excludes a $7$-PD-set contained in this image,
independently of the information set.
\end{proof}

\begin{remark}
Theorem~\ref{thm:counter-fp2} concerns the indicated
stabilizer, not the entire permutation automorphism group
of the $[84,42,\geq40]$ code.  That full code
permutation group has not been determined here, and
no claim about its ability to support full decoding
is deduced from the stabilizer obstruction alone.
\end{remark}

\section{Revised questions and computational complexity}
\label{sec:questions}
The two original conjectures combined a group-size estimate
with the assertion that a group is a full PD-set.  The
counterexamples above show why those assertions must be
separated.  The following is a more precise problem.

\begin{problem}[Partial permutation decoding]
\label{conj:partial}
For the family $C_X(mP_\infty,E_2)$ over $\F_{p^2}$,
with $p\equiv3\pmod4$, determine the largest
integer $s=s(p,m,I)$ for which the geometric stabilizer
$H$ gives an $s$-PD-set relative to an information set $I$.
Compare this value with the largest $s$ obtainable from
the full permutation automorphism group of the code, and
characterize information sets attaining the best values.
\end{problem}

A similar question for the $\F_p$ family asks which values
of $s$ above two are attained by the affine group, as a
function of $p$, $k$, and the chosen information set.
The counting bound \eqref{eq:counting} is a necessary
first filter; its satisfaction does not prove coverage.

It is also important to separate the number of candidate
permutations from the cost of each decoding trial.
If the syndrome is computed by a dense $r\times n$
parity-check matrix, then one trial costs
$O(nr)$ field operations, in addition to $O(n)$
coordinate work, and a worst-case scan of a set $S$
costs $O(|S|nr)$ field operations.  Specialized
syndrome methods can improve this bound, but are not
analyzed here.  The estimates $|\AGL(1,p)|=O(p^2)$
and $|H|=O(p^2)$ describe group cardinalities, not
complete field-operation complexities.  In particular,
these estimates do not establish asymptotic superiority
over standard Reed--Solomon or AG decoding algorithms.

The alphabet sizes $p$ and $p^2$ also vary with the
block length, a fact that must be accounted for in
any bit-complexity comparison.

\section{Summary of the corrected mathematical results}
\begin{enumerate}
\item The $\F_p$ one-point evaluation codes of length $p$
are full-length Reed--Solomon codes, not a new code family
(Theorem~\ref{thm:RS}).
\item The affine permutation group yields a rigorous
positive two-error decoding theorem
(Theorem~\ref{thm:two-pd}).
\item A necessary covering bound excludes the proposed
full-decoding assertions in explicit cases:
$p=37$ over $\F_p$ and $p=7$ over $\F_{p^2}$
(Theorems~\ref{thm:counter-RS} and
\ref{thm:counter-fp2}).
\item An equivariant evaluation code is a submodule of a
permutation module; its Riemann--Roch module need not be
a permutation module even under the original strong degree
hypotheses (Example~\ref{ex:not-permutation}).
\item The main remaining issue is to determine optimal
partial PD-sets, especially for the $\F_{p^2}$-orbit
codes where the construction retains genuinely
hyperelliptic features.
\end{enumerate}

\paragraph{Acknowledgments.}
The original author thanked Bob Guralnick for an orbit
argument and an anonymous referee for helpful comments.
The present AI-aided revision supplies a written proof of
the orbit statement and independently checks the finite
counting counterexamples.  It preserves the original
research motivation and the historical acknowledgments.

\begin{thebibliography}{99}

\bibitem[CLLS]{CLLS}
E. Camps-Moreno, J. B. Lau, H. H. L\'opez, and W. Santos,
``The permutation group of Reed--Solomon codes over arbitrary points,''
\emph{IEEE International Symposium on Information Theory} (2025),
DOI: 10.1109/ISIT63088.2025.11195458;
arXiv:2601.00122.

\bibitem[KW]{KW}
R. Kn\"orr and W. Willems, ``The automorphism groups of generalized
Reed--Muller codes,'' \emph{Ast\'erisque} \textbf{181--182} (1990),
195--207.
  

%\bibitem[Bo1]{Bo1} N. Borne, {\it Une formule de Riemann-Roch
%equivariante pour des courbes,} preprint
%1999
%
%\bibitem[Bo2]{Bo2} -----, {\it Structure de groupe de Grothendieck
%\'equivariant d'une courbe et modules galoisiens,} preprint
%
%\bibitem[Bo3]{Bo3} -----,{\it Modules galoisiens sur les courbes: 
%une introduction,} S\'em. et Congr\`es, SMF \underline{5}(2001)147-159.
%
%\bibitem[B]{B} T. Breuer, {Characters and automorphism
%groups of Riemann surfaces}, London Math. Soc.
%Lecture Notes, 1999.
%
%\bibitem[CR]{CR} C. Curtis, I. Reiner, {Methods of 
%representation theory, I}, Wiley-Interscience, 1981.
% 
%\bibitem[F]{F} W. Fulton, {Algebraic curves},
%Benjamin, 1969.
%
%\bibitem[FH]{FH} ------ and J. Harris, {Representation theory:
%a first course}, Springer-Verlag, 1991.

%\bibitem[GAP]{GAP4}
%The GAP~Group, \emph{GAP -- Groups, Algorithms, and Programming, 
%Version 4.3}; 2002, (\verb+http://www.gap-system.org+).

\bibitem[G]{G} N. G\"ob, ``Computing the automorphism groups of
hyperelliptic function fields,'' available at
\verb+http://front.math.ucdavis.edu/math.NT/0305284+.

\bibitem[GUAVA]{GUAVA}
GUAVA package for GAP 4.4 (\verb+http://www.gap-system.org+),
available at
\verb+http://cadigweb.ew.usna.edu/~wdj/gap/GUAVA/+
or
\newline
\verb+http://www.gap-system.org/Packages/guava.html+

\bibitem[Gu]{Gu} R. Guralnick, private communication, 2004.

\bibitem[H]{H} R. Hartshorne, {\it Algebraic geometry},
Springer-Verlag, 1977.

\bibitem[HvLP]{HvLP}
T. H\o holdt, J. van Lint, R. Pellikaan
``Algebraic-geometric codes,''
in the {\it Handbook of coding theory}, vol 1, pp. 871-961,
(V.S. Pless, W.C. Huffman and R.A. Brualdi Eds.),
Elsevier, Amsterdam 1998.

\bibitem[HP]{HP} W. C. Huffman and V. Pless, 
{\it Fundamentals of 
error-correcting codes}, Cambridge Univ. Press, 2003.

\bibitem[JK1]{JK1} D. Joyner and A. Ksir, 
``Automorphism groups of some AG codes'' available at
\verb+http://front.math.ucdavis.edu/math.AG/0412459+

\bibitem[JK2]{JK2} D. Joyner and A. Ksir, 
``Representations of finite groups on Riemann-Roch
spaces, II'' available at
\verb+http://front.math.ucdavis.edu/math.AG/0312383+

\bibitem[JT]{JT} D. Joyner and W. Traves, 
``Representations of finite groups on Riemann-Roch
spaces'' available at
\verb+http://front.math.ucdavis.edu/math.AG/0210408+

%\bibitem[Ka]{Ka} E. Kani, 
%{\it The Galois-module structure of the 
%space of holomorphic differential forms on a curve,}
%J. reine angew. Math \underline{367}(1986)187-206.
%
%\bibitem[K]{K} B. K\"ock, {\it Galois structure of 
%Zariski cohomology 
%for weakly ramified covers of curves,} 
%math.AG/0207124, available at
%\verb+http://front.math.ucdavis.edu/+
%

\bibitem[L]{L} J. Leon, ``Computing automorphism 
groups of error-correcting codes,'' IEEE Trans. 
Inform. Theory 28(1982)496--511

%\bibitem[L]{L} D. Lorenzini, {An invitation to arithmetic 
%geometry}, Grad. Studies in Math, AMS, 1996.

%\bibitem[MAGMA]{M} 
%W. Bosma, J. Cannon, C. Playoust,
%``The MAGMA algebra system, I: 
%The user language,''
%J. Symb. Comp., \underline{24}(1997)235-265.
%\newline
%(\verb+http://www.maths.usyd.edu.au:8000/u/magma/+).

%\bibitem[MP]{MP} I. Morrison and H. Pinkham,
%{\it Galois Weierstass points and Hurwitz characters},
%Annals of Math., \underline{124}(1986)591-625.
%
%\bibitem[N]{N} S. Nakajima, {\it Galois module structure of
%cohomology groups of an algebraic variety},
%Inv. Math. \underline{75}(1984)1-8
%

\bibitem[S]{S}
C. Schoen, ``On certain modular representations 
in the cohomology 
of algebraic curves,'' J. Alg. \underline{135}(1990)1-18.


%\bibitem[Se]{Se} J.-P. Serre, {Linear representations of 
%finite groups}, Springer-Verlag, 1977.
%
%\bibitem[St]{St} S. Stepanov, {Codes on algebraic curves},
%Kluwer, NY, 1999.

\bibitem[Sti]{Sti} H. Stichtenoth, {\it Algebraic function
fields and codes}, Springer-Verlag, 1993.

\bibitem[TV]{TV}
M. A. Tsfasman and S. G. Vladut, {\it Algebraic-geometric codes},
Mathematics and its Applications, 
Kluwer Academic Publishers, Dordrechet, 1991.

\bibitem[W]{W} S. Wesemeyer, 
``On the automorphism group of 
various Goppa codes,''
IEEE Trans. Info. Theory., 44(1998)630-643.
\end{thebibliography}


\end{document}
